\section*{Hoofdstuk \ref{ch:b3:electronic-spectroscopy} --- Elektronische spectroscopie en fotofysica}

\begin{solution}{exo:b3:electronic-spectroscopy:1}
(a) Verboden volgens de spinregel ($\Delta S = 1$). (b) Verboden volgens de
\omterm{thm:b3:electronic-spectroscopy:laporte}{regel van Laporte}
($g \to g$); zwak zichtbaar via vibronische koppeling. (c) Toegestaan: $B_{1u}$ is de
representatie van $z$ in $D_{2h}$. (d) Verboden door de symmetrie: $A_2$ is
geen van $x$,
$y$, $z$ in $C_{2v}$; de band is zwak.
\end{solution}

\begin{solution}{exo:b3:electronic-spectroscopy:2}
$10^7/349 - 10^7/450 = 28653 - 22222 = \qty{6431}{cm^{-1}}$, \qty{0.80}{eV}.
\end{solution}

\begin{solution}{exo:b3:electronic-spectroscopy:3}
$A = 5700 \times 1 \times 1.0\times10^{-5} = 0.057$; geabsorbeerde fractie $1 - 10^{-0.057} =
0.12$.
\end{solution}

\begin{solution}{exo:b3:electronic-spectroscopy:4}
$k_r = \Phi_F/\tau = \qty{1.5e8}{s^{-1}}$, $\tau_r = \qty{6.7}{ns}$; $k_{nr} = (1 - \Phi_F)/\tau
= \qty{1.0e8}{s^{-1}}$.
\end{solution}

\begin{solution}{exo:b3:electronic-spectroscopy:5}
$\int\varepsilon\,\dd\tilde\nu = 1.0645 \times 5700 \times 5000 = \num{3.03e7}$, dus $f \approx
4.32\times10^{-9} \times 3.03\times10^7 = 0.13$: een toegestane overgang van
gemiddelde sterkte.
\end{solution}

\begin{solution}{exo:b3:electronic-spectroscopy:6}
De poissonfactor $\eu^{-S}S^v/v!$ is maximaal bij $v = \lfloor S\rfloor$. $S = 0.3$: de
0--0-lijn zelf. $S = 1.5$: $v' = 1$, 1.5 keer de 0--0-lijn. $S = 5$: $v' = 4$
en 5 even sterk, 26 keer de 0--0-lijn.
\end{solution}

\begin{solution}{exo:b3:electronic-spectroscopy:7}
$\sum k = \qty{4.2e8}{s^{-1}}$: $\tau = \qty{2.4}{ns}$, $\Phi_F = 0.24$, $\Phi_{\mathrm{ISC}} = 0.71$
(en $\Phi_{\mathrm{IC}} = 0.05$).
\end{solution}

\begin{solution}{exo:b3:electronic-spectroscopy:8}
$\Phi = 0.546 \times 0.46 \times (0.050/0.040) \times (1.4266/1.333)^2 = 0.36$.
\end{solution}

\begin{solution}{exo:b3:electronic-spectroscopy:9}
$\tau_0/\tau = 1.429$ en 1.860: lineair in $[Q]$, dus de levensduur zelf
wordt verkort: \omterm{def:b3:electronic-spectroscopy:quencher}{dynamische uitdoving}. $K_{SV} = 43$ L/mol (beide punten), $k_q = 43/8.0
\times 10^{-9} = \qty{5.4e9}{L.mol^{-1}.s^{-1}}$.
\end{solution}

\begin{solution}{exo:b3:electronic-spectroscopy:10}
\omterm{def:b3:electronic-spectroscopy:quencher}{Statische uitdoving}: de moleculen die nog uitzenden, leven even lang als
voordien; de helft ervan is in de grondtoestand gebonden in een
niet-fluorescerend complex. Met een
associatieconstante $K_a$ is de vrije fractie $1/(1 + K_a[Q])$ en $I_0/I = 1 +
K_a[Q]$, dezelfde vorm als bij Stern-Volmer, maar met onveranderde $\tau$.
\end{solution}

\begin{solution}{exo:b3:electronic-spectroscopy:11}
$\tau_r = \tau/\Phi_F = \qty{14}{ns}$; $\Phi_T = 1 - 0.36 = 0.64$. De gemeten
levensduur wordt verkort door de concurrerende \omterm{def:b3:electronic-spectroscopy:radiationless}{intersystem crossing};
$\tau_r$ is wat emissie alleen zou geven.
\end{solution}

\begin{solution}{exo:b3:electronic-spectroscopy:12}
$\frac12\langle\alpha\beta - \beta\alpha|\alpha\beta + \beta\alpha\rangle = \frac12(1 + 0 - 0 - 1) = 0$, met
$\langle\alpha|\alpha\rangle = \langle\beta|\beta\rangle = 1$, $\langle\alpha|\beta\rangle = 0$ voor elk elektron.
De dipooloperator raakt de spin niet, zodat het overgangsmoment deze
factor nul bevat: $T_1 \leftarrow S_0$ heeft geen intensiteit, tenzij
\omterm{def:b3:many-electron-atoms:spin-orbit}{spin-baankoppeling}
de toestanden mengt.
\end{solution}

\begin{solution}{pb:b3:electronic-spectroscopy:1}
\textbf{1.} $A = 5700 \times 1 \times 1.0\times10^{-4} = 0.57$.
\textbf{2.} $1 - 10^{-0.57} = 0.73$.
\textbf{3.} Een $\varepsilon$ van enkele duizenden: een toegestane overgang
$\pi\to\pi^*$, van gemiddelde sterkte.
\textbf{4.} $1239.8/349 = \qty{3.55}{eV}$.
\textbf{5.} Het absorbeert alleen in het ultraviolet; het zichtbare licht
gaat erdoor.
\textbf{6.} Om het geabsorbeerde licht evenredig te houden met de
concentratie en heropname van het uitgezonden licht te vermijden
(binnenfiltereffecten).
\textbf{7.} $S_0 \to S_1$ (absorptie, daarna \omterm{def:b3:electronic-spectroscopy:radiationless}{vibratierelaxatie}); vanuit $S_1$:
\omterm{def:b3:electronic-spectroscopy:luminescence}{fluorescentie}, \omterm{def:b3:electronic-spectroscopy:radiationless}{interne conversie}, \omterm{def:b3:electronic-spectroscopy:radiationless}{intersystem crossing} naar $T_1$.
\textbf{8.} $k_r = 0.546/19 \times 10^{-9} = \qty{2.87e7}{s^{-1}}$; $\tau_r = \qty{35}{ns}$.
\textbf{9.} $(1 - 0.546)/\tau_0 = \qty{2.39e7}{s^{-1}}$.
\textbf{10.} $28653 - 22222 = \qty{6431}{cm^{-1}}$.
\textbf{11.} De emissie vertrekt vanuit de gerelaxeerde $S_1$ en eindigt op
vibrationeel aangeslagen niveaus van $S_0$, nadat molecuul en oplosmiddel
gerelaxeerd zijn: het foton heeft minder energie, hier genoeg om in het
zichtbare gebied te vallen.
\textbf{12.} $\Phi_F \times (349/450) = 0.42$: 42\,\% van de geabsorbeerde energie.
\textbf{13.} De punten stijgen met 0.594 per \qty{0.010}{mol/L}: een rechte door 1.
\textbf{14.} Kleinste kwadraten door de vijf punten: helling $K_{SV} = \qty{59.4}{L/mol}$,
afsnijding 1.000.
\textbf{15.} $k_q = K_{SV}/\tau_0 = \qty{3.1e9}{L.mol^{-1}.s^{-1}}$.
\textbf{16.} $\tau = \tau_0/(1 + 59.4 \times 0.040) = \qty{5.6}{ns}$.
\textbf{17.} Levensduren meten: bij \omterm{def:b3:electronic-spectroscopy:quencher}{dynamische uitdoving} ligt $\tau_0/\tau$
op dezelfde rechte als $I_0/I$.
\textbf{18.} Tonic bevat geen chloride dat het fluoresceren uitdooft, en is
zuur genoeg om kinine geprotoneerd te houden, zijn fluorescerende vorm.
\textbf{19.} $k_d = 8 \times 8.314 \times 298/(3 \times 0.890\times10^{-3}) =
\qty{7.4e6}{m^3.mol^{-1}.s^{-1}} = \qty{7.4e9}{L.mol^{-1}.s^{-1}}$.
\textbf{20.} $k_q$ is iets minder dan de helft van $k_d$.
\textbf{21.} Ongeveer 42\,\%.
\textbf{22.} Trager: $k_d \propto 1/\eta$, en $k_q$ kan alleen mee dalen.
\textbf{23.} \textbf{$k_q/k_d = 0.42$}: chloride dooft kinine uit met bijna
de helft van de snelheid van diffusiegecontroleerde ontmoetingen.
\end{solution}
